\usepackage{keytheorems}


\usepackage[mathscr]{eucal}% so-called Euler fonts, allows \mathscr
\usepackage{amssymb,amsmath}% allows special arrows like >-> e.g.
\usepackage[utf8]{inputenc}
\usepackage[T1]{fontenc}
\usepackage{lmodern}
\usepackage{amsthm}% allows theorem* , e.g.
\usepackage{comment}% allows multiple lines to be commeted out
\usepackage{graphicx} % Required for inserting images
\usepackage{zref-clever}
\zcsetup{cap,nameinlink}
% Source - https://tex.stackexchange.com/a/150516
% Posted by daleif, modified by community. See post 'Timeline' for change history
% Retrieved 2026-08-04, License - CC BY-SA 3.0

%\usepackage[notref,notcite]{showkeys}
\usepackage{todonotes}
\newcommand{\ttodo}[1]{\todo[color=green!40]{#1}}
\newcommand{\tcomment}[1]{\todo[inline,color=green!40]{#1}}
\newcommand{\dtodo}[1]{\todo[color=Orange!40]{#1}}
\newcommand{\dcomment}[1]{\todo[inline,color=Orange!40]{#1}}

\usepackage{mathtools}
\newcommand\SetSymbol[1][]{\nonscript\:#1\vert\allowbreak\nonscript\:\mathopen{}}
\providecommand\given{} % to make it exist
\DeclarePairedDelimiterX\Set[1]\{\}{\renewcommand\given{\SetSymbol[\delimsize]}#1}

% Save the original command and redefine
\usepackage{tikz-cd}
\usepackage{lmodern}
\usepackage[dvipsnames]{xcolor}

\usepackage{microtype}
\newcommand{\cat}[1]{\mathscr{#1}}%or: \nc{\cat}[1]{\mathcal{#1}}
\usetikzlibrary{arrows}
\usepackage[unicode,colorlinks]{hyperref}
\hypersetup{
  linkcolor=BrickRed,
  citecolor=Green,
  filecolor=Mulberry,
  urlcolor=NavyBlue,
  menucolor=BrickRed,
  runcolor=Mulberry,
  linkbordercolor=BrickRed,
  citebordercolor=Green,
  filebordercolor=Mulberry,
  urlbordercolor=NavyBlue,
  menubordercolor=BrickRed,
  runbordercolor=Mulberry
}
\usepackage{etoolbox,mathtools} %this is needed for the hack removing address indentation
\usepackage{xparse}
\usepackage[all]{xy}
\DeclareMathOperator{\pt}{pt}
%\swapnumbers %pour que le numéro apparaisse devant le théorème
\providecommand\given{}
\newcommand\LocSymbol[1][]{\nonscript\:#1\vert\allowbreak\nonscript\:\mathopen{}}
\DeclarePairedDelimiterX\LocArg[1]{(}{)}{\renewcommand\given{\LocSymbol[\delimsize]}#1}
\newcommand{\Loc}[1]{\operatorname{Loc}\LocArg{#1}}
%\usepackage[a4paper, textheight = 25cm, textwidth = 16cm, top = 2.5cm, centering]{geometry}
\newcommand{\eps}{\varepsilon}

\newcommand{\nc}{\newcommand}
\numberwithin{equation}{section}
\DeclareMathOperator{\im}{im}
\newkeytheorem{Thm}[name=Theorem, sibling=equation, Refname={Theorem,Theorems}]
\newkeytheorem{Lem}[name=Lemma, sibling=Thm, Refname={Lemma,Lemmas}]
\newkeytheorem{Prop}[name=Proposition, sibling=Thm, Refname={Proposition,Propositions}]
\newkeytheorem{Cor}[name=Corollary, sibling=Thm, Refname={Corollary,Corollaries}]
\newkeytheorem{Exa}[name=Example, sibling=Thm, Refname={Example,Examples},style=remark]
\newkeytheorem{Def}[name=Definition, sibling=Thm, Refname={Definition},style=remark]
\newkeytheorem{Rem}[name=Remark, sibling=Thm, Refname={Remark,Remarks},style=remark]
\newkeytheorem{Cons}[name=Construction, sibling=Thm, Refname={Construction},style=remark]
\newkeytheorem{Not}[name=Notation, sibling=Thm, Refname={Notation},style=remark]
\newkeytheorem{Hyp}[name=Hypothesis, sibling=Thm, Refname={Hypothesis},style=remark]
\nc{\dmo}{\DeclareMathOperator}
\dmo{\surj}{surj}
\nc{\bbZ}{\mathbb{Z}}
\nc{\bbN}{\mathbb{N}}

\usepackage[
    backend=biber,
    style=alphabetic,
    maxalphanames=4,
    minalphanames=4,
    maxbibnames=5,
    maxcitenames=5,
    backref=true,
    sortlocale=de_DE,
    natbib=true,
    doi=false,
    isbn=false,
    url=false
]{biblatex}
\addbibresource{bib.bib}
\DeclareMathOperator{\Surj}{Surj}
\DeclareMathOperator{\Pic}{Pic}
\DeclareMathOperator{\Spec}{Spec}
\DeclareMathOperator{\coker}{coker}
\DeclareMathOperator{\Supp}{Supp}
\DeclareMathOperator{\lcm}{lcm}
\DeclareMathOperator{\Epi}{Epi}
\DeclareMathOperator{\Exc}{Exc}
\DeclareMathOperator{\Sp}{Sp}
\DeclareMathOperator{\Hom}{Hom}
\DeclareMathOperator{\Aut}{Aut}
\newcommand{\cO}{\mathcal{O}}
\newcommand{\num}[1]{[#1]}
\makeatletter
\newcommand{\bbQ}{\mathbb{Q}}
\newcommand{\bbS}{\mathbb{S}}
\newcommand{\bbF}{\mathbb{F}}
\newcommand{\pr}{pr}
\newcommand{\gr}{\operatorname{gr}}
\renewcommand*\pod[1]{%
  \allowbreak
  \mathchoice
    {\mkern 18mu}%
    {\mkern  8mu}%
    {\mkern  8mu}%
    {\mkern  8mu}%
  (#1)%
}
\DeclareMathOperator{\Inv}{Inv}
\DeclareMathOperator{\Th}{Th}
\DeclareMathOperator{\Sm}{Sm}
\DeclareMathOperator{\RadId}{RadIdl}
\DeclareMathOperator{\Der}{D}
\DeclareMathOperator{\Idl}{Idl}
\makeatother
\newcommand{\LQ}{\mathop{\mathrm{LQ}}\nolimits}
\usepackage{dsfont}
\newcommand{\unit}{\mathds{1}}
\newcommand{\CAlg}{\operatorname{CAlg}}
\usepackage{bbm}
\DeclareMathOperator{\Mod}{Mod}

\newcommand{\ba}[1]{\mathbb{#1}}
\newcommand{\meet}{\land}
\newcommand{\join}{\lor}
\DeclareMathOperator{\At}{At}
\DeclareMathOperator{\Clop}{Clop}
\newcommand{\R}[1]{\mathfrak{R}(#1)}
\newcommand{\Idem}[1]{\mathrm{IdemIdl}(#1)}
%\DeclareMathOperator{\Idem}{Idl_{\mathrm{idem}}}
\DeclareMathOperator{\Ult}{Ult}
\newcommand{\ideal}[1]{\mathfrak{J}_{#1}}

\newcommand{\BRings}{\mathbf{BoolRings}}
\newcommand{\BAlg}{\mathbf{BoolAlg}}
\DeclareMathOperator{\Spc}{Spc}
\newcommand{\smSpc}{\Spc^{\mathrm{sm}}}


\DeclareMathOperator{\colim}{colim}
\zcRefTypeSetup{equation}{
  Name-sg-ab =,
  name-sg-ab =,
  Name-pl-ab =,
  name-pl-ab =,
  abbrev,
  refbounds = {,(,),},
}
\newkeytheoremstyle{mainthmstyle}{
    headformat=\NAME\ \NUMBER\NOTE
}
\newkeytheorem{MainThm}[
    name=Theorem,
    style=mainthmstyle,
    counter-format=\Alph*,
    Refname={Theorem,Theorems}
]
\setcounter{tocdepth}{1}